\section{Introduction}

Symbolic music generation represents music as structured symbolic events
\citep{huangPopMusicTransformer2020,fradetMiditokPythonPackage2023}, such
as note onsets, durations, velocities, programs, and timing information, rather
than as raw audio waveforms. This representation makes it possible to model
musical structure with sequence models while keeping the output editable as MIDI.
Instead of predicting samples in the audio domain, the system predicts symbolic
events that can later be decoded into notes, instruments, tempo changes, and
other performance-related information.

In this project, MIDI files are converted into discrete token sequences and
modeled with autoregressive neural language models. This follows the same broad
sequence-modeling view used by recent symbolic-music Transformer systems
\citep{hsiaoCompoundWordTransformer2021,ensMMMExploringConditional2020}. The
practical focus is the development of an xLSTM-based system for generating
symbolic music continuations from a given prompt. The work therefore combines
several engineering tasks that
must fit together: selecting a usable MIDI subset, constructing leakage-aware
train/validation/test splits, choosing a token representation, training xLSTM
models, comparing checkpoints, and preparing generated-MIDI evaluation.

The repository\footnote{\url{https://github.com/Alkhatir/duet-of-models}} is
organized for controlled experiments on symbolic music generation using xLSTM
models and, for later thesis comparison, simplified GPT-2-style Transformer
baselines. Although the repository includes infrastructure for both model
families, the current practical phase prioritizes the xLSTM path: tokenizer
selection, scalable xLSTM configurations, training runs, hyperparameter sweeps,
and checkpoint selection.

The main goal of this report is not to claim a final music-generation model, but
to document the current practical state of the pipeline and the experimental
evidence available so far. The report first describes the dataset and split
construction, then explains the preprocessing and REMI+-style tokenization
choices, then presents the xLSTM training setup and model-selection process.
The results and discussion sections summarize what the completed validation
experiments show, what compute was spent, and which conclusions should be
deferred until held-out generation evaluation is finished.
